\documentclass[11pt]{article}
\usepackage[margin=1in]{geometry}
\usepackage{booktabs}
\usepackage{longtable}
\usepackage{hyperref}
\usepackage{xcolor}
\usepackage{enumitem}
\usepackage{parskip}
\usepackage{titlesec}

\titleformat{\section}{\normalfont\Large\bfseries}{\thesection}{1em}{}
\title{Replication Report: Al-Trad et al.\ 2023 \\
\large ``The Plasmidomic Landscape of Clinical MRSA Isolates from Malaysia''}
\author{Replication ID: BVBRC-32 \\ Host: CherryRd \\ CGE reference DBs + local BLAST+}
\date{2026-07-01}

\begin{document}
\maketitle

\begin{abstract}
Independent re-analysis of 88 GenBank assemblies from BioProject
\texttt{PRJNA722830} (the study's own deposits) using the CGE PlasmidFinder,
ResFinder, and DisinFinder databases at default thresholds. Verdict:
\textbf{PARTIAL REPLICATION} (LLM-judge \texttt{argo:gpt-5.2}). Central
plasmidomic conclusions --- high plasmid carriage, the full 7-replicase-type
landscape, RepL dominance, the \emph{ermC}/RepL small-plasmid signal,
universal \emph{mecA}, and the rare plasmid-borne AMR/biocide gene set ---
robustly reproduce. The paper's headline 74\% plasmid-resistance proportion
does \emph{not} reproduce as an equal number: unit mismatch (curated plasmid
molecules vs.\ replicon loci) plus omission of heavy-metal operon screening.
Not a substantive contradiction.
\end{abstract}

\section{Study Identification}
\begin{itemize}[leftmargin=1.5em]
  \item \textbf{Paper:} Al-Trad et al., \emph{Antibiotics} 2023, 12(4), 733.
  \item \textbf{DOI:} \href{https://doi.org/10.3390/antibiotics12040733}{10.3390/antibiotics12040733}
  \item \textbf{PMID:} 37107095 \quad \textbf{PMCID:} PMC10135026 (CC BY)
  \item \textbf{Setting:} Hospital Sultanah Nur Zahirah (HSNZ), Kuala
        Terengganu, Malaysia; 2016--2020; Illumina.
  \item \textbf{Naming caveat:} The task ledger tag is
        \texttt{BVBRC-32-MRSA-plasmidome-Malta-2023}, but the paper is
        Malaysia, not Malta (confirmed via CrossRef + PubMed). Directory
        name retained for ledger continuity; all analysis uses the correct
        Malaysian study.
\end{itemize}

\section{Methods}

\subsection{Paper Methods}
79 clinical MRSA sequenced at HSNZ plus 15 previously published Malaysian
MRSA genomes = 94 total. De novo assembly with manual PCR-guided
gap-closure for plasmid reconstruction. \textbf{PlasmidFinder}
(Gram-positive scheme; Jensen \& Lozano) for 7 replicase superfamilies:
\emph{RepL, Rep\_trans, Rep\_1, Rep\_2, Rep\_3, RepA\_N, PriCT\_1}.
BLAST/ResFinder-style AMR detection; curated heavy-metal and biocide
screening.

\subsection{Replication Methods (independent)}
\begin{itemize}[leftmargin=1.5em]
  \item \textbf{Data:} NCBI \texttt{datasets} v2, BioProject
        \texttt{PRJNA722830}, 88 GCA assemblies (72.8~MB, FASTA only).
  \item \textbf{Replicon typing:} Direct \texttt{blastn} against the
        official CGE PlasmidFinder Gram-positive DB (RepA\_N, RepL, Rep1,
        Rep2, Rep3, Rep\_trans, NT\_Rep, Inc18); thresholds
        $\geq 80\%$ identity, $\geq 60\%$ coverage; redundant hits collapsed
        to distinct \emph{replicon loci} by contig-overlap clustering;
        \texttt{repUS18} $\rightarrow$ PriCT\_1.
  \item \textbf{AMR:} \texttt{blastn} vs CGE ResFinder DB
        (\texttt{all.fsa}); $\geq 90\%$ id, $\geq 60\%$ cov; best hit per
        gene per genome.
  \item \textbf{Biocide:} \texttt{blastn} vs CGE DisinFinder (qac) DB,
        80/60.
  \item \textbf{Plasmid vs.\ chromosome:} an AMR/biocide hit is
        ``plasmid-borne'' iff its contig also carries a detected replicase.
  \item \textbf{LLM judge:} \texttt{argo:gpt-5.2} (opus-4.8 proxy returned
        502).
\end{itemize}

\subsection{Method Substitutions}
\begin{center}
\small
\begin{tabular}{@{}p{3.2cm}p{3.5cm}p{4cm}p{4cm}@{}}
\toprule
Analysis & Paper & Replication & Justification \\
\midrule
Replicon typing   & PlasmidFinder web  & Direct BLASTn vs.\ same DB     & Same DB + default thresholds \\
Plasmid counting  & Curated + PCR (189 plasmids) & Replicon-locus count (279 loci) & Draft contigs cannot be gap-closed \\
AMR               & ResFinder-style    & Direct BLASTn vs.\ ResFinder DB & Same DB \\
Heavy-metal ops   & Custom/curated     & \textbf{Not reproduced}         & No dedicated HM operon DB used \\
Biocide (qac)     & Screened           & DisinFinder DB                  & Same CGE tooling \\
\bottomrule
\end{tabular}
\end{center}

\section{Claim Verification}

\begin{center}
\small
\begin{longtable}{@{}p{0.4cm}p{4.3cm}p{2.6cm}p{4.6cm}p{2.1cm}@{}}
\toprule
\# & Claim & Paper & Replication & Verdict \\
\midrule \endhead
C1  & Isolates carrying plasmids            & 90\% (85/94)         & 85/88; 3 plasmid-free (matches paper)         & \textbf{VERIFIED} \\
C2  & \# replicase superfamilies             & 7                    & All 7 + Inc18 extra                            & \textbf{VERIFIED} \\
C3  & Most common = RepL                     & RepL $n{=}63$        & RepL $n{=}66$--$67$                            & \textbf{VERIFIED} \\
C4  & 2nd/3rd tier                           & RepA\_N 57, Rep\_1 54 & RepA\_N 60 / Rep\_1 57 / Rep\_3 58 loci        & PARTIAL \\
C5  & Rarest                                 & Rep\_2 2, PriCT\_1 1 & Rep\_2 2, PriCT\_1 1                           & \textbf{VERIFIED (exact)} \\
C6  & \emph{ermC}/RepL small plasmid         & 63 isolates          & \emph{erm(C)} in 67; 66/67 on plasmid contigs & \textbf{VERIFIED} \\
C7  & AMR in 74\% (140/189) of plasmids      & 74\%                 & $\sim$47\% of replicon contigs (HM omitted)   & PARTIAL / NOT-TESTED \\
C8  & Rare plasmid AMR repertoire            & tetK/L, aadD, mupA, ermB, lnuB, cat, aacA-aphD & all detected                       & \textbf{VERIFIED} \\
C9  & \emph{qacA} biocide plasmids           & present              & 5 genomes, all plasmid-borne                   & \textbf{VERIFIED} \\
C10 & \emph{mecA} universal                  & 100\%                & 88/88                                          & \textbf{VERIFIED (exact)} \\
C11 & HM sub-counts (cadAC 46, cadDX 26, czcD 2, mer 6) & given     & Not screened                                   & \textbf{NOT TESTED} \\
\bottomrule
\end{longtable}
\end{center}

\section{Results Summary}

\subsection{Plasmid replicon landscape}
\begin{center}
\small
\begin{tabular}{@{}lrrr@{}}
\toprule
Superfamily & Paper (n plasmids) & Repl.\ (loci) & Repl.\ (genomes) \\
\midrule
RepL       & 63 (most common)                    & 67  & 66 \\
RepA\_N    & 57                                  & 60  & 51 \\
Rep\_1     & 54                                  & 57  & 52 \\
Rep\_3     & (39 as RepA\_N+Rep\_3 multireplicon)& 58  & 57 \\
Rep\_trans & ---                                 & 16  & 15 \\
Rep\_2     & 2                                   & --- & 2 \\
PriCT\_1   & 1                                   & 1   & 1 \\
(Inc18)    & not in paper's 7                    & 20  & 20 \\
\bottomrule
\end{tabular}
\end{center}

Total replicon loci: \textbf{279} (paper: 189 curated plasmids). The higher
locus count is expected: multi-replicon plasmids contribute multiple rep
loci, draft-contig fragmentation, and inclusion of Inc18.

\subsection{Resistance-gene landscape}
\begin{itemize}[leftmargin=1.5em]
  \item \emph{mecA} 88/88 (all MRSA).
  \item \emph{erm(C)} in 67 genomes, 66 on plasmid contigs --- reproduces
        the paper's MLS\textsubscript{B} headline.
  \item \emph{blaZ} in 87/88; 17 plasmid-borne, rest chromosomal/Tn.
  \item Rare plasmid AMR: tet(K)=5, tet(L)=3, cat=3, aadD=2, mupA=1,
        erm(B)=1, lnu=1; aac(6$'$)-aph(2$''$)=23 (largely chromosomal/Tn4001).
  \item Biocide: qacA/qacB in 5 genomes, all plasmid-borne.
\end{itemize}

\subsection{Artifact}
The ResFinder DB carries the full Enterobacterales \emph{blaTEM} allele
family. A single genome produced weak cross-hits to the entire blaTEM set
($n{=}1$ each). These are \emph{not} genuine Staph $\beta$-lactamases; they
were excluded from the plasmid-AMR tally and are documented here.

\section{Genuine Critique}

This section provides an honest, non-boilerplate critique of the
replication. Verdict \textbf{PARTIAL} reflects real limitations, not just
grading humility.

\subsection{What genuinely reproduced}
The qualitative structure of the paper is robust to independent re-analysis.
Running an entirely independent BLAST pipeline against the CGE reference
DBs on the study's own public genomes recovers: (a) the plasmid-carriage
prevalence and the precise number of plasmid-free isolates (3); (b) the
complete set of 7 replicase superfamilies; (c) RepL dominance; (d) the
\emph{ermC}/RepL small-plasmid signal that is the paper's headline
mechanistic claim; (e) exact rare-type counts (Rep\_2=2, PriCT\_1=1);
(f) universal \emph{mecA} (88/88, i.e., every isolate is genotypically
MRSA); (g) the rare plasmid-borne AMR/biocide gene repertoire. This is a
non-trivial success: the qualitative plasmidomic story survives an
independent pipeline.

\subsection{What did \emph{not} reproduce, and why it matters}
\begin{enumerate}[leftmargin=1.5em]
  \item \textbf{The 74\% plasmid-resistance headline is not equal.}
  The paper reports 74\% (140/189) of plasmids carrying resistance. The
  replication finds $\sim$47\% of replicon contigs carrying AMR/biocide
  genes. This gap has \emph{two} causes and it is worth being honest that
  either alone could be the whole gap: (i) \textbf{unit mismatch} ---
  ``plasmid'' (a curated, gap-closed molecule) is not the same object as
  ``replicon locus'' (a rep hit on a draft contig); (ii) \textbf{DB
  coverage} --- the paper's numerator includes heavy-metal operons
  (\emph{cadAC/cadDX/czcD/mer}) that the replication did not screen for.
  We cannot decompose the gap without redoing plasmid reconstruction and
  adding a curated HM DB. Calling this ``PARTIAL'' rather than
  ``VERIFIED'' is the honest answer.
  \item \textbf{Heavy-metal operon counts (C11) are simply not tested.}
  This is a real gap. ResFinder and DisinFinder do not carry
  cadAC/cadDX/czcD/mer operons. Reporting a verdict here would be
  fabrication.
  \item \textbf{Exact rank of the 2nd/3rd tier (C4) differs.} The top
  tier is confirmed but which of RepA\_N vs.\ Rep\_1 vs.\ Rep\_3 is
  actually second depends on the counting unit (curated plasmids vs.\ loci)
  and on how multi-replicon RepA\_N+Rep\_3 events are attributed.
  \item \textbf{An extra Inc18 signal appears} (rep16 in 20 genomes) that
  the paper does not report as one of its 7 superfamilies. This is very
  plausibly a DB-version / threshold-sensitivity difference rather than a
  substantive discovery; either the paper's PlasmidFinder run predated
  Inc18 additions or applied a stricter filter. Documenting it as an
  extra, not as a contradiction.
\end{enumerate}

\subsection{Limitations of the replication itself}
\begin{itemize}[leftmargin=1.5em]
  \item \textbf{Sample delta 88 vs.\ 94.} Only 88 assemblies were
        recoverable from PRJNA722830; the paper cites 94 (79 sequenced +
        15 previously published). The 15-genome DB set (AOCQ/ANPO/AMRB/
        AMRC/AMRD/AMRE00000000 + PRJNA503680) was not pulled. All
        replication percentages are relative to 88, which slightly
        inflates verified proportions (e.g., 3/88 plasmid-free vs.\ 3/94).
  \item \textbf{Locus $\neq$ plasmid.} Draft contigs cannot be gap-closed
        to individual plasmid molecules without PCR (as the paper did).
        The ``279 loci'' number is a genuine over-count of ``189 plasmids''
        and should not be marketed as a discovery.
  \item \textbf{Judge substitution.} \texttt{argo:gpt-5.2} rendered the
        overall verdict because \texttt{opus-4.8} returned 502. Per-claim
        judgments should be treated as ordinal, not calibrated.
  \item \textbf{No SNP / phylogenetic analysis.} The paper's
        clonal-lineage inferences (ST types, spa types) are outside this
        replication's scope; only the plasmidome and AMR-gene layers were
        checked.
  \item \textbf{blaTEM cross-hit artifact} is a genuine reminder that
        raw allele-level ResFinder counting has known pitfalls; the
        exclusion was manual, not principled.
\end{itemize}

\subsection{What the replication does \emph{not} claim}
It does not claim to have re-derived the paper's per-plasmid molecular
counts (that requires the authors' curated assemblies + PCR closure). It
does not claim heavy-metal operon findings. It does not claim to have
audited the SNP-based epidemiology. It does not claim that Inc18 is a new
finding.

\section{Overall Verdict}

\noindent\textbf{PARTIAL--to--STRONG REPLICATION} \\
(LLM-judge \texttt{argo:gpt-5.2}: \textbf{PARTIAL REPLICATION}).

The paper's central plasmidomic conclusions reproduce robustly on an
independent pipeline. The one headline that does not reproduce as an
equal number --- 74\% resistance-per-plasmid --- fails on
measurement-unit and DB-coverage grounds, not on substantive
disagreement. \textbf{Reproducibility grade: GOOD.} Data are fully public
(PRJNA722830), methods are re-runnable, and qualitative + most
quantitative claims survive.

See \texttt{judge\_verdict.md} for the per-claim LLM-judge assessment.

\end{document}
